\chapter{Total Parenteral Nutrition}

\section*{Overview}
Total parenteral nutrition (TPN) is the intravenous delivery of complete nutrition, including carbohydrates, protein, lipids, electrolytes, vitamins, and trace elements, bypassing the gastrointestinal tract. It is infused through a central venous catheter when oral or enteral feeding is impossible or insufficient.

\section*{Alternative Names}
TPN, Intravenous hyperalimentation, Parenteral nutrition, Central parenteral nutrition

\section*{Principle}
TPN delivers a sterile, individually compounded nutrient solution directly into the bloodstream, meeting caloric and nutritional needs without enteral absorption. Because the solution is hypertonic, it must be infused through a central venous catheter. The content of macronutrients, electrolytes, and micronutrients is individualized to the patient's metabolic and fluid requirements.

\section*{History}
Total parenteral nutrition was pioneered by Dudrick and Wilmore in the late 1960s, establishing a life-sustaining therapy for patients with intestinal failure.

\section*{Why It Is Done}
TPN is ordered for patients with intestinal failure, short bowel syndrome, prolonged ileus or bowel obstruction, severe pancreatitis, or other conditions in which the gut cannot be used for more than several days. It is also used in severe malnutrition when oral or enteral feeding is not feasible.

\section*{Is This a Primary or Secondary Test or Procedure}
Primary therapy for intestinal failure and a definitive source of nutrition when enteral feeding is not possible.

\section*{How It Is Performed}
A central venous catheter is placed under sterile technique, typically in the subclavian or internal jugular vein or as a PICC line. The pharmacy-compounded TPN solution is infused continuously or cyclically through a dedicated catheter lumen. Glucose, electrolytes, fluid balance, and clinical status are monitored throughout the infusion.

\section*{Preparation}
Nutritional assessment and baseline laboratory studies, including electrolytes and renal and liver function tests, are obtained. Central venous access is inserted with aseptic precautions, and informed consent is obtained.

\section*{Contraindications and Precautions}
TPN should be avoided when the gastrointestinal tract is usable or in patients with hemodynamic instability who cannot tolerate volume loading. Use caution in hyperglycemia, hypertriglyceridemia, severe electrolyte disturbances, or fluid overload, and treat these before initiating therapy. Central line insertion is contraindicated with untreated coagulopathy or local infection.

\section*{Risks}
Catheter-related bloodstream infection, central venous thrombosis, pneumothorax at insertion, hyperglycemia, refeeding syndrome, and electrolyte imbalances. Long-term use can cause hepatic steatosis, cholestasis, and metabolic bone disease.

\section*{Aftercare and Recovery}
Glucose, electrolytes, and liver function are monitored regularly, and the solution is adjusted to maintain metabolic stability. TPN is tapered before discontinuation to avoid hypoglycemia, and oral or enteral feeding resumes as soon as gut function allows.

\section*{Outcome and Follow-up}
Response is assessed through laboratory values such as glucose, electrolytes, liver function tests, triglycerides, and nitrogen balance rather than a single numeric result. The formulation is adjusted to keep these parameters stable.

\section*{Expected Results}
Stable glucose, electrolytes, and fluid balance with adequate nutritional parameters and weight maintenance in the absence of metabolic complications.

\section*{Follow-up}
Periodic metabolic and nutritional laboratory monitoring, with surveillance of liver function and bone mineral density in long-term patients. Patients may be discharged on home TPN with intestinal rehabilitation plans.

\section*{When to Seek Medical Care}
Follow the procedure team's specific instructions. Seek urgent medical assessment for severe or worsening pain, uncontrolled bleeding, fever or signs of infection, trouble breathing, chest pain, fainting, new weakness or confusion, inability to urinate, or any rapidly worsening symptom. The appropriate warning signs depend on the procedure and the patient's medical condition.

\section*{References}
\begin{enumerate}
  \item MedlinePlus. Total parenteral nutrition. U.S. National Library of Medicine; 2025.
  \item Current Medical Diagnosis and Treatment. McGraw-Hill; 2025.
\end{enumerate}

\clearpage
